\documentclass[11pt]{article}
\usepackage[T1]{fontenc}
\usepackage[utf8]{inputenc}
\usepackage{lmodern}
\usepackage[margin=1in]{geometry}
\usepackage{amsmath,amssymb,amsthm,mathtools}
\usepackage[shortlabels]{enumitem}
\usepackage{hyperref}
\input{release_info.tex}

\hypersetup{colorlinks=true,linkcolor=blue,citecolor=blue,urlcolor=blue}

\newtheorem{theorem}{Theorem}[section]
\newtheorem{corollary}[theorem]{Corollary}
\newtheorem{definition}[theorem]{Definition}
\newtheorem{assumption}[theorem]{Assumption}
\newtheorem{proposition}[theorem]{Proposition}
\newtheorem{claim}[theorem]{Claim}
\newtheorem{conjecture}[theorem]{Conjecture}
\theoremstyle{remark}
\newtheorem{remark}[theorem]{Remark}
\theoremstyle{plain}

\title{Reality as a Consensus Protocol:\\
The Fixed-Point Computation That Implements Physics}
\author{B.\ M\"uller \and Kai Xue \and Kale Arnav Anirudha}
\date{\OPHPaperReleaseDate}

\begin{document}
\maketitle
\OPHPaperReleaseBanner

\begin{abstract}
This paper gives the consensus-theoretic core of Observer-Patch Holography (OPH). If physics is built from local observer descriptions that must agree on overlaps, the first question is whether different repair orders lead to different worlds. On the fixed-cutoff collar branch, the local repair step is not left as an abstract rewrite primitive: it is read from exact Markov splice on exact collars or from declared Petz/Fawzi--Renner reconciliation on recoverable collars, and committed only under the declared local-fit contract on the overlaps it touches. That contract makes the inconsistency potential \(\Phi\) a Lyapunov functional for the accepted repair moves, and the fixed-cutoff union-collar gluing package makes overlapping repairs parenthesization-independent on the physical quotient. Under repair completeness, the repair dynamics therefore has a unique schedule-independent normal form. The paper also identifies the loop obstructions to global consistency, shows why physical uniqueness belongs on the gauge quotient and its physical observable algebra rather than on raw microscopic representatives, proves observable-level confluence on the quantum lift even when microscopic representatives differ by gauge relabelings globally or by sector/higher-gauge relabelings inside one declared quotient-local glued state, and gives a fixed-point account of stable operator-algebraic records. This is the finite patch-net theorem package that supports the broader relativity, gauge, particle, and observer branches in the OPH suite.
\end{abstract}

\section{Introduction}

This paper gives a finite patch-net formulation of the OPH consensus picture. A finite graph of
observer patches carries local state spaces and overlap maps. Local recovery-derived repair maps act
asynchronously, each modifying only one patch or a bounded neighborhood, and global consistency
means agreement on every overlap. The central mathematical problem is whether this repair dynamics
admits a unique normal form and how global obstructions to reconciliation are encoded.

The resulting theorem package has two layers. The first layer concerns convergence: the declared
local-fit contract on touched overlaps makes \(\Phi\) a Lyapunov functional for accepted repair
moves, so termination is derived on the finite patch net; the fixed-cutoff union-collar gluing
package then turns competing local repairs into a local diamond on the physical quotient, so
under repair completeness every initial configuration has a unique schedule-independent normal
form. The second layer concerns obstructions: pairwise overlap agreement does not ensure a
global solution, and the obstruction is holonomic. On the abelian branch it is the cycle sum of
edge data; on the genuinely noncentral branch it is a crossed-module \v{C}ech class.

Gauge symmetry enters as invariance under changes of hidden local representation that preserve
overlap data. When the repair step is read only on that overlap-invariant quotient, the
normal-form map descends to the gauge quotient, so physical uniqueness is a quotient statement.
On the quantum lift, the same quotient-local carrier determines a unique terminal state on every
declared physical observable algebra, even when microscopic representative lifts differ by gauge
or sector relabelings inside one quotient-local glued state.
The observation layer is carried by finite observer-accessible record algebras generated by
central or quantitatively stable approximately commuting projectors. These results form the
patch-net formulation used in the broader OPH literature.

This paper proves four core results:

\begin{enumerate}
\item \textbf{Asynchronous confluence.} For the declared accepted repair law, the local-fit contract makes \(\Phi\) a Lyapunov functional and hence gives termination, while the fixed-cutoff union-collar gluing package gives the local diamond on the physical quotient; under repair completeness, every initial state has a unique normal form, independent of update schedule (Theorem~\ref{thm:confluence}).
\item \textbf{Cycle obstruction.} For affine overlap constraints over an abelian group, global consistency holds if and only if the holonomy vanishes on every cycle (Theorem~\ref{thm:holonomy}). The parity triangle gives the minimal frustrated example, and Theorem~\ref{thm:higherdefects} extends the same logic to the crossed-module higher-gauge defect hierarchy used later in OPH.
\item \textbf{Gauge quotient and observable-level confluence.} When local repair is induced on the overlap-invariant quotient, the normal-form map descends to that quotient, and the induced terminal state on every declared physical observable algebra is unique there even when microscopic representatives differ by gauge or sector relabelings inside one quotient-local glued state (Theorems~\ref{thm:gauge}, \ref{thm:obsconfluence}).
\item \textbf{Record algebra and stability.} On the fixed-cutoff observer-accessible surface, central record projectors carry Born/L\"uders measurement directly, and approximate record projectors inherit explicit \((\varepsilon,\delta_{\mathrm{rec}})\) stability bounds on the same event surface (Theorem~\ref{thm:crdt}).
\end{enumerate}

The repair step itself is not a free rewrite primitive here. On the
fixed-cutoff collar branch, a local update is obtained from exact Markov splice or from a
declared Petz/Fawzi--Renner recovery channel and then read on overlap-invariant physical data.
The fixed-point theorem below still isolates the separate remaining burdens cleanly: the declared
repair law already includes the touched-overlap local-fit contract that yields Lyapunov descent of
\(\Phi\) on accepted moves, while the fixed-cutoff gluing package supplies the
parenthesization-independent union-collar glue used for the local diamond. What remains above
that step is repair completeness and, on the Petz branch, the support/CPTP clause stated later in
Proposition~\ref{prop:petz-cptp}.

We also define a fitness functional over the space of candidate reconciliation laws and prove that replicator dynamics monotonically increases mean fitness (Theorem~\ref{thm:replicator}). This gives a clean mathematical model for the selection of physical law.

The results here are exact theorems about a computational model. The companion OPH manuscript develops the broader physics branches conditionally and keeps structural theorems, scaling-limit branches, calibration-sector outputs, and phenomenological continuations distinct~\cite{OPH}.


\section{Patch Nets, Overlaps, and Global Consistency}

\begin{definition}[Patch net]\label{def:patchnet}
Let $G=(V,E)$ be a finite connected graph. Each vertex $i\in V$ is an \textbf{observer patch} with finite local state space $S_i$. The global state space is
\[
\Sigma := \prod_{i\in V} S_i.
\]
For each edge $e=\{i,j\}\in E$, let $I_e$ be an interface alphabet and let
\[
\pi_{i,e}:S_i\to I_e,
\qquad
\pi_{j,e}:S_j\to I_e
\]
be the interface projection maps. A global state $s=(s_i)_{i\in V}\in\Sigma$ is \textbf{consistent on edge $e=\{i,j\}$} iff
\[
\pi_{i,e}(s_i)=\pi_{j,e}(s_j).
\]
The global consistency set is
\[
C:=\bigl\{s\in\Sigma:\forall\, e=\{i,j\}\in E,\ \pi_{i,e}(s_i)=\pi_{j,e}(s_j)\bigr\}.
\]
\end{definition}

For exposition we use a finite pairwise-overlap graph. The hypergraph version is straightforward: replace edges by hyperedges and pairwise equality by a common interface label on each hyperedge. Nothing in the proofs depends on the pairwise restriction.

The picture: each observer holds a local state, and neighboring observers share an interface through which they can compare notes. A universe-state is physically admissible exactly when all neighbors agree on their shared data. This is a constraint satisfaction problem (CSP), and the consistent states are the codewords.

\begin{definition}[Inconsistency potential]\label{def:potential}
For each edge $e$, choose a weight $w_e>0$ and a function $d_e:I_e\times I_e\to\mathbb{R}_{\ge 0}$ with $d_e(a,b)=0 \iff a=b$. On the declared fixed-cutoff branch, \(d_e\) is the overlap score used by the local acceptance contract on that interface. Define
\[
\Phi(s)
:=
\sum_{e=\{i,j\}\in E}
w_e\, d_e\!\bigl(\pi_{i,e}(s_i),\pi_{j,e}(s_j)\bigr).
\]
Then $s\in C \iff \Phi(s)=0$.
\end{definition}

So $\Phi$ is the total disagreement energy of the universe. Consistent states have zero energy. Everything else is frustrated.


\section{Asynchronous Reconciliation and the Main Fixed-Point Theorem}

\begin{definition}[Recovery-derived local repair law]\label{def:repair}
Fix for each patch \(i\) a finite collar chart \(A_i\!-\!B_i\!-\!D_i\) around the overlaps touched
by \(i\), together with a fixed local decoder from repaired collar data back to the finite patch
label at \(i\). A \textbf{law} \(\lambda\) is a family of local repair maps
\[
T_i^\lambda:\Sigma\to\Sigma
\qquad (i\in V)
\]
such that \(T_i^\lambda\) changes only the state of patch \(i\) (or, more generally, only a
bounded neighborhood of \(i\)), and the local update is induced by one of the declared OPH
recovery moves on that collar, where \(\omega_{A_iB_i}(s)\) denotes the \(A_i\cup B_i\) marginal
of the current collar state encoded by \(s\):
\begin{enumerate}[label=(\roman*),leftmargin=2em]
\item exact Markov splice on the collar, using Theorem~\ref{thm:splice} when
      \(I(A_i:D_i\mid B_i)=0\); or
\item a declared recoverability channel
\[
(\mathrm{id}_{A_i}\otimes \mathcal R_i)(\omega_{A_iB_i}(s)),
\qquad
\mathcal R_i=\mathcal R_{\sigma_i,\mathcal N_i},
\]
      with \(\mathcal R_i\) in the Petz/Fawzi--Renner class of
      Definition~\ref{def:oph-repair-map} and Theorem~\ref{thm:splice}.
\end{enumerate}
The decoder back to \(S_i\) is bookkeeping for the finite patch presentation; the physical content
is the repaired collar state on the declared fixed-cutoff branch. Write \(s\to_i t\) iff
\(t=T_i^\lambda(s)\neq s\). Let \(\to := \bigcup_{i\in V}\to_i\), and let \(\to^*\) be its
reflexive-transitive closure. A state \(s\in\Sigma\) is a \textbf{normal form} iff
\(T_i^\lambda(s)=s\) for all \(i\in V\).
\end{definition}

\begin{definition}[Touched-overlap potential and accepted local-fit contract]\label{def:localfit}
For each repair site \(i\), let
\[
E_i^{\mathrm{touch}}
:=
\bigl\{
e\in E:\text{the interface data on }e\text{ may change under }T_i^\lambda
\bigr\}.
\]
Define the touched-overlap potential
\[
\Phi_i(s)
:=
\sum_{e=\{u,v\}\in E_i^{\mathrm{touch}}}
w_e\, d_e\!\bigl(\pi_{u,e}(s_u),\pi_{v,e}(s_v)\bigr).
\]
On the declared fixed-cutoff branch, a recovery-derived candidate is committed only if it
strictly lowers this touched-overlap score:
\[
s\to_i t \implies \Phi_i(t)<\Phi_i(s).
\]
This is the patch-net form of the regulator-side monotone local-fit contract carried by the
declared repair package.
\end{definition}

\begin{definition}[Overlap-associative union-collar gluing]\label{def:assoc}
Fix \(s\in\Sigma\) and two enabled repairs \(s\to_i t\), \(s\to_j u\). Write
\[
E_{ij}^{\mathrm{touch}}
:=
E_i^{\mathrm{touch}}\cup E_j^{\mathrm{touch}}.
\]
The declared repair branch is \textbf{overlap-associative} if the following hold.
\begin{enumerate}[label=(\roman*),leftmargin=2em]
\item If \(E_i^{\mathrm{touch}}\cap E_j^{\mathrm{touch}}=\varnothing\), the two accepted local
updates have disjoint support on the declared branch and therefore commute.
\item If \(E_i^{\mathrm{touch}}\cap E_j^{\mathrm{touch}}\neq\varnothing\), there is a finite union
collar \(U_{ij}\) covering the interfaces in \(E_{ij}^{\mathrm{touch}}\) such that the physical
glued state on \(U_{ij}\) is parenthesization-independent on the quotient, in the sense of
Proposition~\ref{prop:assocglue}, and the local decoders/lifts of
Definition~\ref{def:repair} are restriction-compatible on nested collars.
\end{enumerate}
This is the concrete compatibility package used below to complete the local diamond.
\end{definition}

\begin{remark}[Inputs and remaining burdens]\label{rem:repair-burdens}
The repair step is therefore not an abstract rewrite primitive. Its declared inputs are the
fixed-cutoff collar chart, either exact Markov splice or a chosen Petz/Fawzi--Renner recovery
channel, a local decoder/lift back to the finite patch presentation, and the touched-overlap
local-fit contract of Definition~\ref{def:localfit}, together with the support-local
disjoint-commutation clause and the restriction-compatible union-collar package of
Definition~\ref{def:assoc}. The
parenthesization-independent quotient-local glue used there is supplied by
Proposition~\ref{prop:assocglue} from the fixed-cutoff center-sector / higher-gauge gluing
package. What the present theorem package still assumes rather than derives is repair
completeness. On the Petz branch, full CPTP action on all inputs also requires the support clause
recorded later in Proposition~\ref{prop:petz-cptp}. On broader branches one must also prove that
the declared union-collar compatibility is preserved under refinement or branch change.
\end{remark}

The theorem package separates the imported repair-law data from the remaining theorem-local inputs cleanly:

\begin{assumption}[Repair completeness]\label{ass:complete}
$s\in C \iff \forall\, i\in V,\; T_i^\lambda(s)=s.$
\end{assumption}

Normal forms are exactly the globally consistent states. The dynamics is neither too weak (missing some inconsistencies) nor too strong (repairing things that were fine).

\begin{proposition}[Accepted repair moves are Lyapunov-decreasing]\label{prop:lyapunov-descent}
For the accepted repair law of Definitions~\ref{def:repair} and \ref{def:localfit}, every enabled
repair strictly decreases the inconsistency potential:
\[
s\to t \implies \Phi(t)<\Phi(s).
\]
\end{proposition}

\begin{proof}
Write \(s\to_i t\). If \(e\notin E_i^{\mathrm{touch}}\), then the interface data on \(e\) are
unchanged by \(T_i^\lambda\), so the corresponding term in \(\Phi\) is the same for \(s\) and
\(t\). Therefore
\[
\Phi(t)-\Phi(s)=\Phi_i(t)-\Phi_i(s)<0
\]
by Definition~\ref{def:localfit}.
\end{proof}

\begin{proposition}[Termination from the OPH Lyapunov functional]\label{prop:lyapunov-termination}
Under Proposition~\ref{prop:lyapunov-descent}, every repair sequence is finite; equivalently, the
repair relation \(\to\) is terminating.
\end{proposition}

\begin{proof}
Because each \(S_i\) is finite, the global state space \(\Sigma=\prod_{i\in V}S_i\) is finite, so
the value set \(\Phi(\Sigma)\subset\mathbb{R}_{\ge 0}\) is finite as well. Along any nontrivial
repair step \(s\to t\), Proposition~\ref{prop:lyapunov-descent} gives \(\Phi(t)<\Phi(s)\). An infinite repair
sequence would therefore produce an infinite strictly descending chain in the finite ordered set
\(\Phi(\Sigma)\), which is impossible.
\end{proof}

\begin{proposition}[Local confluence from overlap-associative gluing]\label{prop:local-confluence}
For the accepted repair law of Definitions~\ref{def:repair}, \ref{def:localfit}, and
\ref{def:assoc}, the repair relation is locally confluent.
\end{proposition}

\begin{proof}
Take \(s\to_i t\) and \(s\to_j u\).
If \(E_i^{\mathrm{touch}}\cap E_j^{\mathrm{touch}}=\varnothing\), then
Definition~\ref{def:assoc}(i) gives
\[
T_j^\lambda(T_i^\lambda(s))=T_i^\lambda(T_j^\lambda(s)).
\]
So with \(v:=T_j^\lambda(T_i^\lambda(s))\) we have \(t\to_j v\) and \(u\to_i v\).

Assume now \(E_i^{\mathrm{touch}}\cap E_j^{\mathrm{touch}}\neq\varnothing\). Let \(U_{ij}\) be the
union collar from Definition~\ref{def:assoc}(ii), and let \(\overline v_{ij}(s)\) denote the
quotient-local glued state on \(U_{ij}\) determined by the exterior marginals of \(s\).
Proposition~\ref{prop:assocglue} makes this state independent of whether the local
splice/recovery is parenthesized as \(i\) then \(j\) or \(j\) then \(i\). Because the local
decoders are restriction-compatible on nested collars, the one-site repairs \(s\to_i t\) and
\(s\to_j u\) are precisely the \(i\)- and \(j\)-restrictions of \(\overline v_{ij}(s)\).
Applying that same restriction-compatible decoder to the unfinished subcollar therefore produces
the complementary accepted local step from \(t\) and from \(u\) into representatives of the same
quotient-local union-collar state. Hence there exists \(v\in\Sigma\) with \(t\to^* v\) and
\(u\to^* v\), where \(v\) is any representative of \(\overline v_{ij}(s)\).
\end{proof}

\begin{theorem}[Asynchronous confluence / fixed-point law]\label{thm:confluence}
For the accepted repair law of Definitions~\ref{def:repair}, \ref{def:localfit}, and
\ref{def:assoc}, under Assumption~\ref{ass:complete}, every initial state $s\in\Sigma$ has a unique normal form
\[
\operatorname{nf}_\lambda(s)\in C,
\]
and every maximal asynchronous repair execution from $s$ terminates at that same state. The terminal state is independent of update order.
\end{theorem}

\begin{proof}
By Proposition~\ref{prop:lyapunov-termination}, the repair relation \(\to\) is terminating.
By Proposition~\ref{prop:local-confluence}, it is locally confluent. Newman's
lemma~\cite{newman42} therefore makes it confluent. A terminating confluent repair relation has a
unique normal form reachable from each initial state. By Assumption~\ref{ass:complete}, the
normal forms are exactly \(C\). Every maximal execution terminates and reaches
\(\operatorname{nf}_\lambda(s)\), and that terminal state is independent of update order.
\end{proof}

\begin{corollary}[Objective law is schedule-independent]\label{cor:objective}
Let $M:\Sigma\to Y$ be any observable. Under the hypotheses of Theorem~\ref{thm:confluence}, $M(\operatorname{nf}_\lambda(s))$ is independent of the asynchronous update schedule. If physical law is identified with the map $s\mapsto M(\operatorname{nf}_\lambda(s))$, then physical law is objective.
\end{corollary}

\begin{proof}
All schedules from the same initial $s$ terminate at $\operatorname{nf}_\lambda(s)$, so all yield the same $M$-value.
\end{proof}

At this stage objectivity is identified with schedule-independent convergence of the repair
dynamics; Theorem~\ref{thm:obsconfluence} sharpens this to schedule-independent convergence of the
physical observable algebra even when microscopic representatives are not unique.


\section{Why Local Agreement Is Not Enough: Cycle Holonomy and Frustration}

Now we show that the story has a twist: just because every pair of neighbors agrees does not mean the whole system is consistent.

\begin{theorem}[Cycle-obstruction / holonomy criterion]\label{thm:holonomy}
Let $A$ be an abelian group, and let $G=(V,E)$ be a connected graph with an arbitrary orientation on each edge. For each oriented edge $e:u\to v$, assign a label $b_e\in A$. Consider the affine consistency equations
\[
x_v - x_u = b_e
\qquad
\text{for every oriented edge } e:u\to v,
\]
where $x_v\in A$ are unknown patch labels. A global solution $x:V\to A$ exists if and only if for every cycle $C\subseteq G$,
\[
\sum_{e\in C}\varepsilon_C(e)\, b_e = 0,
\]
where $\varepsilon_C(e)=+1$ if the cycle traverses $e$ in the chosen orientation and $-1$ otherwise.
\end{theorem}

\begin{proof}
\textbf{Necessity.} Suppose $x$ is a solution. Summing the edge equations around any cycle $C$,
\[
\sum_{e:u\to v\in C}\varepsilon_C(e)\,(x_v-x_u)
=
\sum_{e\in C}\varepsilon_C(e)\,b_e.
\]
The left side telescopes to $0$ because every vertex appears once with $+$ sign and once with $-$ sign.

\textbf{Sufficiency.} Fix a root $r\in V$. For any vertex $v$, choose a path $P_{r\to v}$ and define
\[
x_v := \sum_{e\in P_{r\to v}} \varepsilon_{P_{r\to v}}(e)\, b_e,
\qquad x_r:=0.
\]
If $P$ and $P'$ are two paths from $r$ to $v$, traversing $P$ followed by the reverse of $P'$ yields a cycle. By the vanishing-holonomy assumption, the total signed sum is zero, so $x_v$ is well-defined. For any edge $e:u\to v$, extending a path to $u$ by that edge gives $x_v = x_u + b_e$.
\end{proof}

\begin{corollary}[Parity triangle: pairwise consistency is not enough]\label{cor:triangle}
Take $A=\mathbb{Z}_2$ on the triangle $A$-$B$-$C$-$A$ with edge labels $b_{AB}=0$, $b_{BC}=0$, $b_{CA}=1$. Each individual edge equation is satisfiable. But the global system is not: the cycle sum is $0\oplus 0\oplus 1 = 1\neq 0$.
\end{corollary}

This example shows that pairwise consistency does not imply a global solution. The obstruction is
carried by the cycle.

\begin{corollary}[Stable defects as frustrated holonomy]\label{cor:defects}
Define the defect energy
\[
\Phi_b(x) := \sum_{e:u\to v\in E} w_e\,\mathbf{1}\!\left[x_v - x_u \neq b_e\right].
\]
If the cycle-holonomy condition fails, then $\min_{x:V\to A}\Phi_b(x)>0$. Every minimizer contains irreducible residual inconsistency.
\end{corollary}

\begin{proof}
If $\min_x\Phi_b(x)=0$, some assignment satisfies all edge equations, contradicting Theorem~\ref{thm:holonomy}.
\end{proof}

Residual inconsistencies of this type cannot be removed by local repair moves. In the OPH
interpretation they are stable topological defects of the reconciliation dynamics.

\begin{theorem}[Higher-gauge defect hierarchy]\label{thm:higherdefects}
Let a finite overlap nerve carry crossed-module defect data \((g_{ij},h_{ijk})\) for a compact crossed module
\[
H\xrightarrow{\partial} G.
\]
Under local rechartings by
\[
C^1(N,H)\rtimes C^0(N,G),
\]
the nonabelian \v{C}ech class
\[
q=[(g,h)]\in \check H^2(N,H\to G)
\]
is invariant. Strict global reconciliation exists if and only if \(q=0\), and nonzero \(q\) labels stable fixed-cutoff higher-gauge defects.
\end{theorem}

\begin{proof}
The allowed rechartings are exactly the crossed-module coboundaries. So they preserve the class and strictify the weak gluing data precisely when the class vanishes. A nonzero class therefore gives a topologically protected residual obstruction, in the full crossed-module defect hierarchy rather than only in the abelian truncation.
\end{proof}


\section{Gauge Symmetry as Implementation Hiding}

\begin{definition}[Gauge action]\label{def:gauge}
Let $\Gamma = \prod_{i\in V}\Gamma_i$ act on $\Sigma$ componentwise. The action is a \textbf{gauge action} if for every $e=\{i,j\}\in E$,
\[
\pi_{i,e}(\gamma_i\cdot x)=\pi_{i,e}(x)
\qquad
\forall\, x\in S_i,\ \forall\, \gamma_i\in \Gamma_i.
\]
Gauge changes alter hidden local representations but do not alter overlap data.
\end{definition}

Gauge transformations change hidden local representations while leaving overlap data fixed. Write
\[
q:\Sigma\to \Sigma/\Gamma,
\qquad
q(s)=[s],
\]
for the gauge-orbit map. A \textbf{physical repair law} is the family of quotient-local maps
\[
\overline T_i^\lambda:\Sigma/\Gamma\to \Sigma/\Gamma
\]
induced by the recovery-derived collar updates of Definition~\ref{def:repair}. A
\textbf{representative repair family} is any choice of lifts
\[
T_i^\lambda:\Sigma\to\Sigma
\]
such that
\[
q\circ T_i^\lambda
=
\overline T_i^\lambda\circ q.
\]
This is the finite patch-net form of saying that the repair step is defined first on
overlap-invariant physical data and only then lifted to hidden representatives. The remaining
problem here is therefore not to invent a repair rule, but to establish repair completeness on
the declared fixed-cutoff branch and the support/CPTP clause on the Petz branch where that
channel is used. The touched-overlap local-fit contract gives Lyapunov \(\Phi\)-descent on
accepted moves, and Proposition~\ref{prop:assocglue} together with
Definition~\ref{def:assoc} supplies the quotient-local compatibility package used for the local
diamond. Stability of that package under refinement or branch change is a separate question.
The point here is also that no extra gauge-covariance axiom is needed once repair is formulated
on the quotient.

\begin{theorem}[Gauge quotient theorem]\label{thm:gauge}
Under the gauge action of Definition~\ref{def:gauge}, any representative lift of a physical repair
law as just defined, and the hypotheses of Theorem~\ref{thm:confluence},
\[
q\bigl(\operatorname{nf}_\lambda(\gamma\cdot s)\bigr)
=
q\bigl(\operatorname{nf}_\lambda(s)\bigr)
\qquad
\forall\, \gamma\in\Gamma,\ \forall\, s\in\Sigma.
\]
Hence the normal-form map descends to the quotient:
\[
\overline{\operatorname{nf}}_\lambda:\Sigma/\Gamma \to \Sigma/\Gamma,
\qquad
[s]\mapsto [\operatorname{nf}_\lambda(s)].
\]
\end{theorem}

\begin{proof}
Suppose $s\to_i t$, so $t=T_i^\lambda(s)\neq s$. If $s'=\gamma\cdot s$, then
\[
q\bigl(T_i^\lambda(s')\bigr)
=
\overline T_i^\lambda\bigl(q(s')\bigr)
=
\overline T_i^\lambda\bigl(q(s)\bigr)
=
q\bigl(T_i^\lambda(s)\bigr).
\]
Thus gauge-equivalent inputs induce the same repaired orbit, and by induction the orbit reached
after any repair sequence depends only on the initial orbit. Every maximal repair sequence from
$s$ ends at $\operatorname{nf}_\lambda(s)$ by Theorem~\ref{thm:confluence}, so the terminal orbit
$q(\operatorname{nf}_\lambda(s))$ depends only on $q(s)=[s]$. This makes
\[
[s]\longmapsto [\operatorname{nf}_\lambda(s)]
\]
well-defined on $\Sigma/\Gamma$.
\end{proof}

\begin{corollary}[Gauge-invariant law]\label{cor:gaugeinv}
If $M:\Sigma\to Y$ is gauge-invariant ($M(\gamma\cdot s)=M(s)$ for all $\gamma$), then $M(\operatorname{nf}_\lambda(s))$ depends only on the gauge orbit $[s]$, not the representative.
\end{corollary}

\begin{proof}
Because $M$ is gauge-invariant, it factors through the orbit map: $M=\overline M\circ q$ for some
$\overline M:\Sigma/\Gamma\to Y$. Theorem~\ref{thm:gauge} gives
\[
q\bigl(\operatorname{nf}_\lambda(\gamma\cdot s)\bigr)
=
q\bigl(\operatorname{nf}_\lambda(s)\bigr),
\]
hence
\[
M\bigl(\operatorname{nf}_\lambda(\gamma\cdot s)\bigr)
=
\overline M\!\left(q\bigl(\operatorname{nf}_\lambda(\gamma\cdot s)\bigr)\right)
=
\overline M\!\left(q\bigl(\operatorname{nf}_\lambda(s)\bigr)\right)
=
M\bigl(\operatorname{nf}_\lambda(s)\bigr).
\]
\end{proof}

\begin{definition}[Physical observable algebra on the quantum lift]\label{def:physobs}
Fix a finite patch region or union collar \(R\) on the declared fixed-cutoff quantum lift of
Appendix~\ref{app:consensus-fixed-cutoff-port}. Its \textbf{physical observable algebra}
\(\mathcal A_{\mathrm{phys}}(R)\) is the fixed-point collar algebra under the compact boundary
redundancy action on the ordinary or central-defect branch, or the corresponding quotient-local
algebra on the genuinely noncentral branch. The central sector projectors carried by the collar
decomposition belong to \(Z(\mathcal A_{\mathrm{phys}}(R))\). A \textbf{physical observable} is
any \(X\in \mathcal A_{\mathrm{phys}}(R)\). For a microscopic representative \(s\in\Sigma\), let
\(\omega_R^s\) denote the induced state on \(\mathcal A_{\mathrm{phys}}(R)\).
\end{definition}

\begin{theorem}[Observable-level confluence on the quantum lift]\label{thm:obsconfluence}
Under the hypotheses of Theorems~\ref{thm:confluence} and \ref{thm:gauge}, every initial orbit
\([s]\in\Sigma/\Gamma\) has a unique quotient normal form
\[
\overline{\operatorname{nf}}_\lambda([s])\in q(C).
\]
Fix a declared region \(R\). Let \(t,u\in\Sigma\) be terminal microscopic representatives reached
by representative lifts of maximal repair sequences from initial states in the same orbit
\([s]\). If the induced \(R\)-collar states of \(t\) and \(u\) are representatives of the same
quotient-local glued state in the sense of
Proposition~\ref{prop:assocglue}, then for every physical observable
\(X\in\mathcal A_{\mathrm{phys}}(R)\),
\[
\omega_R^t(X)=\omega_R^u(X).
\]
Hence all physical observables converge to the same overlap-consistent values even when the
microscopic terminal representatives differ by gauge relabelings globally or by sector/higher-gauge
relabelings inside one declared quotient-local glued state.
\end{theorem}

\begin{proof}
Theorems~\ref{thm:confluence} and \ref{thm:gauge} give the unique quotient normal form
\(\overline{\operatorname{nf}}_\lambda([s])=[\operatorname{nf}_\lambda(s)]\in q(C)\). Now let
\(t\) and \(u\) be as stated. By Corollary~\ref{cor:obsglue}, representatives of the same
quotient-local glued state induce the same state on the physical observable algebra
\(\mathcal A_{\mathrm{phys}}(R)\), including the central sector projectors carried by that
algebra. Therefore \(\omega_R^t(X)=\omega_R^u(X)\) for every
\(X\in\mathcal A_{\mathrm{phys}}(R)\).
\end{proof}

\begin{remark}[Inputs and remaining boundary for observable-level confluence]\label{rem:obsconfluence}
Theorem~\ref{thm:obsconfluence} does not add a new repair hypothesis beyond the D1 package. Its
inputs are exactly the representative-level confluence theorem, quotient descent of the repair
law, Definition~\ref{def:physobs}, and Corollary~\ref{cor:obsglue} from the fixed-cutoff
union-collar gluing package. What remains open is not fixed-cutoff observable uniqueness on that
carrier, but extension of the same statement to broader refinement-stable branches where the
declared union-collar compatibility is only approximate or where the chosen physical observable
algebra itself changes under refinement.
\end{remark}

\begin{corollary}[Inert ancillary refinement does not change physical law]\label{cor:ancilla}
Let \(K=\prod_i K_i\) be a finite ancillary state space and define \(\Sigma^\eta:=\Sigma\times K\). Lift the repair maps by
\[
T_i^{\lambda,\eta}(s,k):=(T_i^\lambda(s),k).
\]
If \(M^\eta:\Sigma^\eta\to Y\) ignores the ancillary factor,
\[
M^\eta(s,k)=M(s),
\]
with \(M\) gauge-invariant on \(\Sigma\), then
\[
M^\eta(\operatorname{nf}_\lambda^\eta(s,k))
=
M(\operatorname{nf}_\lambda(s))
\]
for all \((s,k)\in\Sigma^\eta\).
\end{corollary}

\begin{proof}
Because the ancillary factor is inert,
\[
\operatorname{nf}_\lambda^\eta(s,k)=\bigl(\operatorname{nf}_\lambda(s),k\bigr).
\]
Therefore \(M^\eta(\operatorname{nf}_\lambda^\eta(s,k))=M(\operatorname{nf}_\lambda(s))\), and Corollary~\ref{cor:gaugeinv} supplies gauge-orbit independence.
\end{proof}

Physical uniqueness therefore holds on the quotient by gauge or implementation hiding. The same
statement remains unchanged under inert ancillary stabilization.


\section{Record Algebras and the Operator Observation Layer}

\paragraph{Inputs used here.}
From the fixed-cutoff collar package we use only the observer-accessible finite-dimensional algebra
on one completed compare/write/verify slice, the declared pointer and overlap-sector projectors
read on that same slice, and the trace-distance control on restored accessible states when
restoration is invoked. No continuum lift and no broader observer-metaphysical premise is used
here.

\begin{definition}[Exact record algebras and approximate record presentations]\label{def:lattice}
Fix one completed observer-accessible slice at cycle \(t\), and let
\(\mathcal A_t^{\mathrm{acc}}\) be the corresponding finite-dimensional accessible algebra. An
\emph{exact record presentation} is a family of orthogonal projectors
\(\{\widehat P_a(t)\}_a\subset Z(\mathcal A_t^{\mathrm{acc}})\) whose generated algebra
\[
\mathcal Z_{\mathrm{rec}}(t)
:=
\mathrm{Alg}\bigl(\{\widehat P_a(t)\}_a\bigr)
\]
is finite and commutative.

An \emph{approximate record presentation} on the same declared readout slots is a family of
projectors \(\{P_a(t)\}_a\subset \mathcal A_t^{\mathrm{acc}}\) together with an exact record
presentation \(\{\widehat P_a(t)\}_a\) such that
\[
\delta_{\mathrm{rec}}(t)
:=
\max_a \|P_a(t)-\widehat P_a(t)\|
\]
is finite.
\end{definition}

\begin{theorem}[Record algebra, Born-L\"uders update, and quantitative stability]\label{thm:crdt}
Fix one completed observer-accessible slice at cycle \(t\).
\begin{enumerate}
\item The exact record projectors generate a finite commutative central algebra
\(\mathcal Z_{\mathrm{rec}}(t)\subset Z(\mathcal A_t^{\mathrm{acc}})\).
\item For every event \(E\) in the finite event algebra generated by \(\mathcal Z_{\mathrm{rec}}(t)\),
\[
\mathbb P_t(E)=\operatorname{Tr}\!\bigl(\rho_t \widehat P_E(t)\bigr),
\qquad
\rho_t\!\mid_E
=
\frac{\widehat P_E(t)\rho_t \widehat P_E(t)}
{\operatorname{Tr}(\rho_t \widehat P_E(t))}.
\]
\item If no accepted repair between \(t\) and \(t+1\) touches the support of \(\widehat P_E(t)\),
then the next read of \(E\) has probability \(1\).
\item If \(\{P_a(t)\}_a\) is an approximate record presentation with modulus
\(\delta_{\mathrm{rec}}(t)\), then
\[
\|[P_a(t),P_b(t)]\|\le 4\,\delta_{\mathrm{rec}}(t)
\]
for all declared record projectors \(P_a(t),P_b(t)\). Moreover, for every declared elementary
record event \(a\) and every restored accessible state \(\widetilde\rho_t\) satisfying
\[
\|\widetilde\rho_t-\rho_t\|_1\le \varepsilon,
\]
one has
\[
\Bigl|
\operatorname{Tr}\!\bigl(\widetilde\rho_t P_a(t)\bigr)
-
\operatorname{Tr}\!\bigl(\rho_t \widehat P_a(t)\bigr)
\Bigr|
\le
\varepsilon+\delta_{\mathrm{rec}}(t).
\]
\end{enumerate}
\end{theorem}

\begin{proof}
Because the exact record projectors lie in the center of \(\mathcal A_t^{\mathrm{acc}}\), they
generate a finite commutative central algebra. Finite-dimensional projective measurement on that
commuting algebra gives the Born trace and the L\"uders conditioned state, proving (1) and (2).

If no accepted repair touches the support of \(\widehat P_E(t)\), then the next completed read is
performed on the same central projector surface. After conditioning on \(E\), the state lies in
the range of \(\widehat P_E(t)\), so the next read of the same event has probability \(1\). This
gives (3).

For (4), write
\[
[P_a(t),P_b(t)]
=
[P_a(t)-\widehat P_a(t),P_b(t)]
+
[\widehat P_a(t),P_b(t)-\widehat P_b(t)],
\]
because \([\widehat P_a(t),\widehat P_b(t)]=0\). Since every projector has operator norm at most
\(1\),
\[
\|[P_a(t),P_b(t)]\|
\le
2\|P_a(t)-\widehat P_a(t)\|
+
2\|P_b(t)-\widehat P_b(t)\|
\le
4\,\delta_{\mathrm{rec}}(t).
\]
For the probability bound,
\[
\Bigl|
\operatorname{Tr}\!\bigl(\widetilde\rho_t P_a(t)\bigr)
-
\operatorname{Tr}\!\bigl(\rho_t \widehat P_a(t)\bigr)
\Bigr|
\le
\Bigl|\operatorname{Tr}\!\bigl((\widetilde\rho_t-\rho_t)P_a(t)\bigr)\Bigr|
+
\Bigl|\operatorname{Tr}\!\bigl(\rho_t(P_a(t)-\widehat P_a(t))\bigr)\Bigr|.
\]
The first term is bounded by \(\|\widetilde\rho_t-\rho_t\|_1\|P_a(t)\|\le \varepsilon\), and the
second by \(\|\rho_t\|_1\|P_a(t)-\widehat P_a(t)\|\le \delta_{\mathrm{rec}}(t)\). Hence the total
error is at most \(\varepsilon+\delta_{\mathrm{rec}}(t)\).
\end{proof}

\paragraph{Merge boundary.}
What is closed here is the fixed-cutoff operator-algebraic observation surface: a central record
algebra on the exact readout slice, Born/L\"uders measurement on its event projectors, exact
repeated-read stability under the untouched-support hypothesis, and explicit
\((\varepsilon,\delta_{\mathrm{rec}})\) control when practical readout projectors are only close
to that central surface. What remains beyond this theorem is not another finite-cutoff record
lemma but the broader export problem: proving on richer branches that physically relevant pointer
surfaces stay close to one such central record algebra and that the same control survives
refinement and continuation.


\section{Law-Space Selection and Observer Emergence}

This section studies a simple meta-selection model on law space. The aim is to formalize one
criterion for favoring schedule-robust, observer-supporting, and simple laws; the replicator
dynamics below is part of the model, not a claim about literal cosmological dynamics.

We begin by defining what it means for a law to support observers. An observer is treated
operationally as a \textbf{persistent predictive module}: a subgraph that maintains a stable
record algebra and uses its output law to predict its boundary's future behavior.

\begin{definition}[Schedule robustness]\label{def:robustness}
Fix distributions $\mu$ over initial conditions and $\nu$ over asynchronous schedules, and a gauge-invariant observable $M$. For a law $\lambda$, define
\[
\mathcal{R}_M(\lambda)
:=
\Pr_{s\sim\mu,\;\sigma,\tau\sim\nu}
\!\left[
M\bigl(\operatorname{nf}^{\sigma}_\lambda(s)\bigr)
=
M\bigl(\operatorname{nf}^{\tau}_\lambda(s)\bigr)
\right].
\]
If Theorem~\ref{thm:confluence} holds for $\lambda$, then $\mathcal{R}_M(\lambda)=1$.
\end{definition}

\begin{definition}[Observer yield]\label{def:yield}
Let $X_t^\lambda$ denote the stationary process obtained by repeated local perturbation plus
reconciliation under law $\lambda$. For each subgraph $U\subseteq V$, let
\(\mathcal Z_U^{\mathrm{rec}}(t)\) be the declared exact record algebra on its observer-accessible
surface, or the reference exact algebra when only an approximate record presentation is available,
and let $Y_U(t)$ be the corresponding finite outcome variable induced by that record algebra. Then $U$ is
\((\eta,\varepsilon,h)\)-observer-like if it is record-stable:
\[
d_{\mathrm{TV}}\!\bigl(\operatorname{Law}(Y_U(t+h)),\operatorname{Law}(Y_U(t))\bigr)\le \eta,
\]
and predictive:
\[
I\bigl(Y_U(t);\; X_{\partial U,\,t+1:t+h}^\lambda\bigr)\ge\varepsilon.
\]
Define
\[
\mathcal{O}_{\eta,\varepsilon,h}(\lambda)
:=
\mathbb{E}\!\left[
\#\left\{
U\subseteq V:
U \text{ is } (\eta,\varepsilon,h)\text{-observer-like}
\right\}
\right].
\]
\end{definition}

\begin{definition}[Law fitness]\label{def:fitness}
Let $K(\lambda)$ be a description-length penalty. Define
\[
f(\lambda) = \alpha\,\mathcal{R}_M(\lambda) + \beta\,\mathcal{O}_{\eta,\varepsilon,h}(\lambda) - \gamma\,K(\lambda),
\]
with $\alpha,\beta,\gamma>0$.
\end{definition}

\begin{theorem}[Replicator monotonicity on law space]\label{thm:replicator}
Let $\Lambda=\{\lambda_1,\dots,\lambda_m\}$ be candidate laws with population weights $x_i(t)$ under replicator dynamics:
\[
\dot{x}_i = x_i(f_i - \bar{f}),
\qquad
f_i := f(\lambda_i),
\qquad
\bar{f} := \sum_{j=1}^m x_j f_j.
\]
Then
\[
\frac{d}{dt}\bar{f} = \operatorname{Var}_x(f) \ge 0.
\]
Mean fitness is nondecreasing, and strictly increasing unless all extant laws have the same fitness.
\end{theorem}

\begin{proof}
Direct computation:
\[
\frac{d}{dt}\bar{f}
=
\sum_i \dot{x}_i f_i
=
\sum_i x_i(f_i - \bar{f})f_i
=
\sum_i x_i f_i^2 - \bar{f}^2
=
\operatorname{Var}_x(f) \ge 0.
\]
Equality iff all $f_i$ on the support of $x$ are equal.
\end{proof}

This theorem records the monotonicity property of the meta-selection model.


\section{Connection to Observer-Patch Holography}

The formalism above is the computational skeleton of Observer-Patch Holography (OPH). The observer patches carry von Neumann algebras on a holographic screen $S^2$, the overlap projections are restrictions to shared subalgebras, and the consistency condition is algebraic state agreement on overlaps.

The bridge to physics works as follows in the broader paper stack:

\begin{itemize}
\item The patch net becomes a net of subregion algebras on $S^2$.
\item Overlap Consistency, one of the canonical OPH axioms, is the algebraic version of Definition~\ref{def:patchnet}.
\item The recoverability clause of the canonical Recoverable Generalized Entropy axiom provides the collar factorization $\rho_{ABD} = \bigoplus_\alpha p_\alpha\,\rho^{(\alpha)}_{Ab_L^\alpha}\otimes\rho^{(\alpha)}_{b_R^\alpha D}$ and the declared Petz/Fawzi--Renner recovery channels from which the local repair moves are built (see Appendix~\ref{app:splice} and Definition~\ref{def:oph-repair-map}).
\item Gauge symmetry as implementation hiding (Theorem~\ref{thm:gauge}) becomes the edge-sector fusion structure from which OPH conditionally reconstructs a compact gauge group and, under the MAR admissibility package, selects the Standard Model gauge group $\mathrm{SU}(3)\times\mathrm{SU}(2)\times\mathrm{U}(1)/\mathbb{Z}_6$.
\item Stable defects (Corollary~\ref{cor:defects}) become the topologically protected excitations that OPH identifies with particles.
\item The record-algebra theorem (Theorem~\ref{thm:crdt}) provides the formal basis for the fixed-cutoff observation layer in OPH, where records are carried by central or quantitatively stable approximately commuting projectors in overlap centers.
\end{itemize}

The companion manuscripts develop a conditional gravity branch from entanglement equilibrium and modular geometry, the SM gauge-group closure from edge-sector admissibility plus MAR, and, on the string/worldsheet lane, a fixed-cutoff heat-kernel edge-sector result on the declared compact-gauge branch together with a worldsheet-like reorganization only if a controlled large-$N_{\mathrm{edge}}$ regime distinct from the physical color rank $N_c=3$ exists; any further critical-superstring lift remains conjectural. Input-dependent cosmological statements and later phenomenological continuations are kept explicitly separate from that recovered-core claim set.

This paper provides the finite patch-net foundation for the broader companion suite. When that
suite is described using labels such as \(D3\)--\(D5\) or \(D7\)--\(D9\), those are internal
documentary node labels from the companion derivation ledger~\cite{OPH}; they are not external
references, and that companion ledger is the place to look up the definitions.


\section{Discussion and Open Problems}

This paper proves the fixed-point consensus spine of OPH: schedule-independent normal forms,
holonomy obstructions, gauge-quotient invariance, and the fixed-cutoff operator-record theorem.
It also gives a
clean law-selection meta-model. The conditional relativity chain and the realized Standard Model
structural chain remain the recovered core. The capacity relation is separate and input-dependent.
Downstream phenomenology requires additional premises beyond the consensus results proved here.

Six directions stand out as immediate next steps:

\begin{enumerate}
\item \textbf{Computational complexity.} What is the complexity of computing $\operatorname{nf}_\lambda(s)$? The finite-valued Lyapunov termination bound from the touched-overlap descent of \(\Phi\) is a worst-case measure, but typical instances may converge much faster. Is there a polynomial bound for natural classes of patch nets?

\item \textbf{Coarse-graining.} When does reconciliation commute with renormalization? If you coarse-grain the patch net and then reconcile, do you get the same result as reconciling first and then coarse-graining? The conditions under which these operations commute would connect the discrete model to continuum field theory.

\item \textbf{Defect classification and refinement-limit transportability.} The fixed-cutoff hierarchy extends from abelian frustrations to crossed-module classes \(q\in \check H^2(N,H\to G)\). The remaining task is to connect those higher-gauge defect sectors to the refinement-stable transportable sector category used in the broader compact-gauge reconstruction lane.

\item \textbf{Observable-level confluence beyond the declared fixed-cutoff physical algebra.} Theorem~\ref{thm:obsconfluence} closes the fixed-cutoff quantum-lift statement when microscopic representatives differ by gauge relabelings globally or by sector/higher-gauge relabelings on the same declared quotient-local glued state. The remaining question is whether an analogous observable theorem survives on broader refinement-stable branches where the union-collar compatibility is only approximate or where the physical observable algebra itself changes under refinement.

\item \textbf{Recovery-channel control.} The fixed-cutoff collar branch already expresses local repair through exact splice or Petz/Fawzi--Renner recovery, commits candidate moves only under the touched-overlap local-fit contract that yields Lyapunov descent of \(\Phi\), and derives the local diamond from support-local disjoint commutation together with parenthesization-invariant union-collar gluing on the physical quotient. What remains is to prove on natural refinement-stable branches that these accepted recovery moves are repair-complete for overlap consistency, preserve support-local disjoint commutation and nested-collar restriction compatibility, including the support/CPTP clause on the Petz branch, and keep the union-collar package stable under refinement or branch changes.

\item \textbf{Scaling-limit bridge to the downstream gravity/gauge stack.} Under what additional regularity conditions does the consensus protocol emit the prime geometric cap pair and support ordered cut-pair rigidity on that extracted limit, so that the D3--D5 gravity chain and the D7--D9 gauge/matter chain become available? This is also the point at which the broader Phase-II and Phase-III branches must remain disciplined: dark-sector, baryogenesis, spectroscopy, and string/worldsheet topics require extra premises beyond what this paper itself proves.
\end{enumerate}

The last item is the bridge to the downstream gravity/gauge branches. Failure of that bridge would
leave the fixed-point theorems proved here unchanged while revising the downstream gravity, gauge,
or continuation sectors.

\subsection{Conditional BFT and QECC Extensions}
\label{subsec:conditional-bft-qecc}

The consensus formalism of \textrm{OPH} has natural analogies to classical and quantum
distributed Byzantine agreement. Observer patches correspond to protocol nodes, overlap repair
corresponds to a quorum vote, and the repair fixed-point corresponds to a consensus state. Under
explicit structural assumptions (quorum size $\geq 2f+1$, partial synchrony, one-vote-per-view,
certificate semantics, and \textrm{DLS}-style view-change), a \textrm{QBFT}-style interpretation
of \textrm{OPH} repair satisfies safety and liveness (Appendix~\ref{app:bft-qecc},
Theorem~\ref{thm:qbft-safety}). On the fixed-cutoff collar branch used here, the repair map is
already written in exact-splice / Petz form; what remains conditional is the \textrm{CPTP}
property on all inputs, which requires either full-rank $\mathcal{N}(\sigma)$ or an explicit
domain restriction, and trace-preserving completion is not automatic when
$\mathcal{N}(\sigma)$ has a non-trivial kernel
(Proposition~\ref{prop:petz-cptp}). A quantum error-correcting interpretation is possible under
additional topological structure, but the code-distance / min-cut equality requires explicit
conditions on logical-operator homology and boundary geometry (Claim~\ref{claim:qecc-mincut}). All
of these extensions are conditional or conjectural and are not part of the core theorem package of
Paper~4.


\appendix
\section{Quantum/Algebraic Lift: Markov-Collar Splice Theorem}\label{app:splice}

This appendix records the algebraic splice statement used to relate the finite patch-net model to
the OPH collar formalism.

An observer is written as
\[
O=(P,\mathcal{A}(P),\rho,R),
\]
where $P$ is the screen patch, $\mathcal{A}(P)$ the local von Neumann algebra, $\rho$ the local state, and $R$ the record algebra.

\begin{theorem}[Markov-collar splice theorem, exact and controlled]\label{thm:splice}
Suppose a collar tripartition $A$-$B$-$D$ has exact Markov decomposition
\[
\rho_{ABD}
=
\bigoplus_{\alpha}
p_\alpha\,
\rho^{(\alpha)}_{A b_L^\alpha}
\otimes
\rho^{(\alpha)}_{b_R^\alpha D}.
\]
Let $\sigma_{b_R^\alpha D'}^{(\alpha)}$ be any family of normalized environment states compatible with the same right-boundary sectors. Define
\[
\rho'_{AB D'}
=
\bigoplus_{\alpha}
p_\alpha\,
\rho^{(\alpha)}_{A b_L^\alpha}
\otimes
\sigma_{b_R^\alpha D'}^{(\alpha)}.
\]
Then for every observable $X$ supported on $A\cup b_L$,
\[
\operatorname{Tr}(X\rho'_{AB D'})
=
\operatorname{Tr}(X\rho_{ABD}).
\]
Now fix one finite-dimensional collar model and let
\[
\mathfrak M_{A:B:D}
:=
\left\{
\tau_{ABD}: I(A:D\mid B)_\tau=0
\right\},
\]
with exact-Markov distance modulus
\[
\delta^{\mathrm M}_{A:B:D}(\varepsilon)
:=
\sup\left\{
\inf_{\tau\in\mathfrak M_{A:B:D}}\|\omega-\tau\|_1:
I(A:D\mid B)_\omega\le\varepsilon
\right\}.
\]
Then
\[
\delta^{\mathrm M}_{A:B:D}(\varepsilon)\to0
\qquad
(\varepsilon\downarrow0).
\]
Hence if \(I(A:D\mid B)_\omega\le \varepsilon\) and \(\widetilde\omega_\varepsilon\in\mathfrak M_{A:B:D}\) is chosen so that
\[
\|\omega-\widetilde\omega_\varepsilon\|_1
\le
\delta^{\mathrm M}_{A:B:D}(\varepsilon),
\]
the corresponding exact splice \(\widetilde\omega'_\varepsilon\) satisfies
\[
\left|
\operatorname{Tr}(X\omega)-\operatorname{Tr}(X\widetilde\omega'_\varepsilon)
\right|
\le
\|X\|_\infty\,
\delta^{\mathrm M}_{A:B:D}(\varepsilon)
\]
for every observable \(X\) supported on \(A\cup b_L\).

Independently, if \(I(A:D\mid B)_\omega\le \varepsilon\), then there exists a recovery map \(\mathcal R_{B\to BD}\) such that
\[
\left\|
\omega_{ABD}
-
(\mathrm{id}_A\otimes \mathcal R_{B\to BD})(\omega_{AB})
\right\|_1
\le
2\sqrt{1-e^{-\varepsilon}}
\le
2\sqrt{\varepsilon}.
\]
\end{theorem}

\begin{proof}
The exact splice statement is the usual blockwise factorization argument:
\[
\operatorname{Tr}(X\rho'_{AB D'})
=
\sum_\alpha
p_\alpha\,
\operatorname{Tr}\!\left(
X\, \rho^{(\alpha)}_{A b_L^\alpha}
\right)
\operatorname{Tr}\!\left(\sigma_{b_R^\alpha D'}^{(\alpha)}\right).
\]
Each right factor is normalized, so the value agrees with the same computation for \(\rho_{ABD}\).

For the controlled statement, compactness of the fixed finite-dimensional state space and continuity of conditional mutual information imply \(\delta^{\mathrm M}_{A:B:D}(\varepsilon)\to0\): otherwise one could find a sequence with \(I(A:D\mid B)\to0\) staying a fixed trace distance away from every exact Markov state, contradicting convergence of a subsequence to an exact Markov limit point. Once \(\widetilde\omega_\varepsilon\) is chosen, the exact splice identity for \(\widetilde\omega_\varepsilon\) gives
\[
\left|
\operatorname{Tr}(X\omega)-\operatorname{Tr}(X\widetilde\omega'_\varepsilon)
\right|
=
\left|
\operatorname{Tr}\!\left[X(\omega-\widetilde\omega_\varepsilon)\right]
\right|
\le
\|X\|_\infty\,
\delta^{\mathrm M}_{A:B:D}(\varepsilon).
\]
The final inequality is the standard Fawzi--Renner recovery bound~\cite{fawziRenner15}.
\end{proof}

This appendix therefore uses exact splice identities in only two regimes: literal exact Markovity, or a controlled collar family on one fixed finite-dimensional model for which \(\delta^{\mathrm M}_{A:B:D}(\varepsilon)\to0\). Small one-shot conditional mutual information is not silently upgraded to an exact normal form.

\input{tex_fragments/CONSENSUS_TECHNICAL_SUPPLEMENT_PORT.tex}

\section{Conditional Distributed-Systems and QECC Extensions of the Consensus Formalism}\label{app:bft-qecc}
\input{appendix_B_bft_qecc_extensions}

\begin{thebibliography}{99}

\bibitem{OPH}
B.~M\"uller, K.~Xue, and P.~Nguyen,
\emph{Observers Are All You Need},
2026.\\
Available at \url{https://oph-book.floatingpragma.io/papers/observers_are_all_you_need.pdf}.

\bibitem{newman42}
M.~H.~A.~Newman,
``On theories with a combinatorial definition of `equivalence',''
\emph{Ann.\ of Math.} \textbf{43} (1942), no.~2, 223--243.

\bibitem{FLP85}
M.~J.~Fischer, N.~A.~Lynch, and M.~S.~Paterson,
``Impossibility of distributed consensus with one faulty process,''
\emph{J.\ ACM} \textbf{32} (1985), no.~2, 374--382.

\bibitem{Lamport1982}
L.~Lamport, R.~Shostak, and M.~Pease,
``The Byzantine generals problem,''
\emph{ACM Trans.\ Program.\ Lang.\ Syst.} \textbf{4} (1982), no.~3, 382--401.

\bibitem{DLS88}
C.~Dwork, N.~A.~Lynch, and L.~Stockmeyer,
``Consensus in the presence of partial synchrony,''
\emph{J.\ ACM} \textbf{35} (1988), no.~2, 288--323.

\bibitem{Moniz2020}
H.~Moniz,
\emph{The Istanbul BFT Consensus Algorithm},
arXiv:2002.03613, 2020.

\bibitem{Saltini}
R.~Saltini et al.,
\emph{QBFT Formal Specification and Verification}.
Available at \url{https://github.com/Consensys/qbft-formal-spec-and-verification}.

\bibitem{Petz1986}
D.~Petz,
``Sufficient subalgebras and the relative entropy of states of a von Neumann algebra,''
\emph{Commun.\ Math.\ Phys.} \textbf{105} (1986), no.~1, 123--131.

\bibitem{FagnolaUmanita2010}
F.~Fagnola and V.~Umanit\`a,
``Generators of detailed balance quantum Markov semigroups,''
\emph{Infinite Dimensional Analysis, Quantum Probability and Related Topics}
\textbf{13} (2010), no.~3, 459--486.

\bibitem{Junge2018}
M.~Junge et al.,
``Universal recovery maps and approximate sufficiency of quantum relative entropies,''
\emph{Ann.\ Henri Poincar\'e} \textbf{19} (2018), no.~8, 2505--2555.

\bibitem{KL97}
E.~Knill and R.~Laflamme,
``Theory of quantum error-correcting codes,''
\emph{Phys.\ Rev.\ A} \textbf{55} (1997), no.~2, 900--911.

\bibitem{Kitaev2003}
A.~Kitaev,
``Fault-tolerant quantum computation by anyons,''
\emph{Ann.\ Phys.} \textbf{303} (2003), no.~1, 2--30.

\bibitem{Dennis2002}
E.~Dennis, A.~Kitaev, A.~Landahl, and J.~Preskill,
``Topological quantum memory,''
\emph{J.\ Math.\ Phys.} \textbf{43} (2002), no.~9, 4452--4505.

\bibitem{BCW98}
H.~Buhrman, R.~Cleve, and A.~Wigderson,
``Quantum vs.\ classical communication and computation,''
in \emph{Proceedings of STOC 1998}, pp.~63--68.

\bibitem{fawziRenner15}
O.~Fawzi and R.~Renner,
``Quantum conditional mutual information and approximate Markov chains,''
\emph{Commun.\ Math.\ Phys.} \textbf{340} (2015), 575--611,
arXiv:1410.0664.

\end{thebibliography}

\end{document}
